% ============================================================
%  Section 1: Introduction
% ============================================================

\section{Introduction}
\label{sec:introduction}

Natural Language Processing~(NLP) has undergone a profound transformation over the
past decade, evolving from hand-crafted rule-based systems through statistical models
to deeply learned neural representations.  Central to this progression has been the
development of \emph{pre-trained language models} that transfer broad linguistic
knowledge to specialised downstream tasks with minimal task-specific supervision.

Among these advances, \bert{}~\cite{devlin2019bert} stands out as one of the most
influential contributions to modern NLP.  Released by Devlin~et~al.\ at Google in
2018, \bert{} introduced a fundamentally different paradigm by pre-training a deep
\emph{bidirectional} Transformer encoder on massive unlabelled corpora and subsequently
fine-tuning the resulting representations for a wide array of tasks.

\subsection{Motivation and Context}

Prior to \bert{}, the dominant pre-training strategies were inherently
\emph{unidirectional}: models such as ELMo~\cite{peters2018deep} used
shallow concatenations of independently trained left-to-right and right-to-left
language models, while the Generative Pre-Training model~(GPT)~\cite{radford2018improving}
trained a strictly left-to-right Transformer.  These designs imposed a critical
limitation: each token's representation was conditioned only on tokens to its left
(or only to its right), preventing the model from jointly exploiting both preceding
and succeeding context.

\bert{} resolves this limitation by employing a \emph{Masked Language Model}~(MLM)
objective that randomly obscures a subset of input tokens and requires the model to
predict them from their \emph{full} bilateral context.  This single design choice
yields substantially richer contextual representations.

\subsection{Limitations of Preceding Approaches}

The three principal shortcomings addressed by \bert{} are summarised below.

\begin{enumerate}
  \item \textbf{Inability to capture long-range dependencies.}  Recurrent
        architectures such as Long Short-Term Memory~(LSTM) networks suffer from
        vanishing gradients, limiting their capacity to encode relationships between
        tokens that are far apart in a sequence.

  \item \textbf{Shallow contextual word embeddings.}  Static embedding methods
        (e.g., Word2Vec, GloVe) assign a single fixed vector to each word type,
        regardless of context.  Polysemous words such as \emph{bank} are therefore
        represented identically in ``river bank'' and ``central bank.''

  \item \textbf{Limited cross-task transferability.}  Prior models were commonly
        re-trained from scratch for each new task, discarding knowledge that could
        be shared across related problems.
\end{enumerate}

\subsection{Contributions of This Report}

This report provides a self-contained technical exposition of \bert{}.  Specifically,
it:
\begin{itemize}
  \item describes the Transformer encoder architecture and its hierarchical
        representation learning capabilities~(\cref{sec:architecture});
  \item derives and interprets the scaled dot-product self-attention
        mechanism~(\cref{sec:attention});
  \item details the token, segment, and positional embedding
        scheme~(\cref{sec:input});
  \item formalises the Masked Language Modelling and Next Sentence Prediction
        pre-training objectives~(\cref{sec:pretraining});
  \item outlines the fine-tuning procedure for representative downstream
        tasks~(\cref{sec:finetuning}); and
  \item critically evaluates \bert{}'s contributions, advantages, and
        limitations~(\cref{sec:discussion}).
\end{itemize}
